\section{Empirical Evaluation}
\label{sec:experiments}

The empirical study is designed as a behavioral characterization of
\textsc{QuaRK}, rather than as a search for a single favorable benchmark.  Its
main component uses controlled synthetic processes for which the temporal
mechanism, task functional, dependence strength, ambient dimension, and
measurement budget can be varied independently.  Two smaller components then
probe transfer beyond that controlled setting: deterministic chaotic systems,
which are deliberately outside the stochastic assumptions of the theory, and
three real-world time-series regression datasets used only as external sanity
checks.  The theory-aligned generators are defined in
\cref{app:dependent-families}, the controlled task functionals in
\cref{app:functionals}, the chaotic stress tests in
\cref{app:chaotic-stress}, and the real-world sanity checks in
\cref{app:real-world-sanity}.

\subsection{Empirical questions}
\label{sec:empirical-questions}

\paragraph{Q1: Which temporal computations can the final reservoir features
support?}
We move from one-step prediction to distributed linear memory, smooth
second-order memory, delayed recall, and a sparse cross-lag interaction.  The
resulting hierarchy separates current or near-future information from long
memory and nonlinear temporal composition.

\paragraph{Q2: How do task memory and process persistence interact with the
reservoir contraction?}
The delay, memory-decay, and process-persistence parameters are varied
independently.  This tests whether the preferred contraction setting changes
when the data or target requires longer retained history, and whether excessive
memory can become detrimental.

\paragraph{Q3: Are the observed strengths and failures stable across distinct
dependence mechanisms?}
Stable vector VARMA, finite-state hidden Markov observations, and Gaussian
GARCH(1,1) provide linear recurrence, latent regime persistence, and
conditional heteroskedasticity.  Shared tasks permit direct comparison across
families, while family-specific targets test filtering and volatility-state
recovery.

\paragraph{Q4: What happens when the ambient input dimension grows while the
quantum input width remains fixed?}
A dedicated study increases the ambient dimension under both sparse-signal and
distributed-signal constructions.  It reports prediction, projection
distortion, classical projection cost, and width sensitivity.  This tests the
width-control interpretation of \cref{prop:jl}; it does not claim end-to-end
runtime independence from the original dimension.

\paragraph{Q5: When does \textsc{QuaRK} improve on, match, or underperform
classical representations?}
The principal controls are an echo-state representation, a matched-dimensional
random nonlinear map, and a classical JL representation, each followed by the
same Mat\'ern readout and comparable validation budget.  Raw-window ridge or
Mat\'ern regression is retained as a low-cost reference when computationally
tractable.

\paragraph{Q6: Which resource choices produce useful operating points?}
Encoded width, reservoir count, observable locality, and contraction are varied
separately.  Predictive error is reported together with feature dimension,
logical width, reservoir executions, state preparations, and temporal channel
applications, so that accuracy gains can be interpreted against their cost.

\paragraph{Q7: How sensitive is the method to finite measurement budgets?}
Exact features are compared with local-Pauli classical-shadow estimates under
both a frozen-readout protocol and a noisy-feature training protocol.  The
former connects directly to \cref{cor:measurement-stability}; the latter is a
practical stress test outside that fixed-readout result.

\paragraph{Q8: Does the architecture remain useful outside the theorem's data
assumptions, and do the controlled findings survive a limited external check?}
Lorenz--63 and Mackey--Glass trajectories test forecasting under sensitive,
nonlinear dynamics without asserting beta mixing.  Three heterogeneous
real-world TSER datasets provide a small sanity check rather than a claim of
broad real-world validation.

\subsection{Common protocol and reporting principles}
\label{sec:empirical-protocol}

\paragraph{Three empirical tiers.}
The primary tier consists of stationary beta-mixing synthetic families and
bounded targets covered by the analytical assumptions.  The second tier uses
chaotic dynamical systems whose empirical behavior is explicitly separated
from the theorem.  The third tier contains three real-world datasets for which
stationarity and mixing are not assumed or estimated.  Conclusions are labeled
according to the tier from which they arise.

\paragraph{Controlled data and task suites.}
Unless an experiment states otherwise, the synthetic protocol uses five
predeclared data-generating seeds, window length \(w=25\), stride \(s=100\),
\(N_{\mathrm{tr}}=5000\) chronological training windows, and
\(N_{\mathrm{te}}=1000\) chronological test windows.  The exact family
parameters and burn-in rules are listed in
\cref{app:dependent-family-protocol}; task directions, normalizations, delays,
and family-specific targets are fixed in
\cref{app:functional-allocation}.  Every theory-aligned synthetic label lies in
\([-1,1]\).

\paragraph{Splits and leakage control.}
Hyperparameters are selected only within the training portion using a fixed
inner chronological training--validation split.  The selected candidate is
refitted on the complete training portion before the test set is evaluated.
True forecasting labels may use future observations, but the corresponding
model input always ends at the stated prediction origin.  Latent HMM states and
GARCH variances are used only to construct synthetic teacher targets.  The
chaotic and real-world protocols use the chronological or official splits
specified in \cref{app:chaotic-protocol,app:real-world-protocol}.

\paragraph{Predeclared finite candidate families.}
For each main comparison, the full candidate set is enumerated before
validation scores are inspected.  A candidate may specify projection and
reservoir seeds, \(n\), \(R\), \(k\), \(\lambda_r\), the observable family,
and Mat\'ern parameters \((\xi,\nu,\eta)\).  Candidates covered by
\cref{cor:finite-family-selection} share the same window length and stride and
use a finite grid with \(\nu>1\).  The final report records the grid, its
cardinality \(L\), and the selected candidate.  Continuous or adaptively
expanded searches are labeled as outside the finite-family corollary.

\paragraph{Classical controls.}
The three principal controls are: (i) an echo-state network with representation
dimension matched to \(p=R|\cO|\); (ii) a matched-dimensional random nonlinear
feature map; and (iii) a classical \(d\to n\) JL map.  All three use the same
Mat\'ern KRR readout and readout-selection grid as \textsc{QuaRK}.  Raw-window
ridge and raw-window Mat\'ern KRR are reported where feasible, but are not
described as dimension- or computation-matched controls.

\paragraph{Readout fitting and metrics.}
An implementation solving
\((K+\lambda_{\mathrm K}I)\bm\alpha=\bm y\) uses
\(\lambda_{\mathrm K}=N_{\mathrm{tr}}\eta\) under
\cref{eq:krr-objective}.  Controlled tasks report MSE and
\[
  \mathrm{NRMSE}_{\mathrm{tr}}
  :=\frac{\sqrt{\mathrm{MSE}}}{\widehat\sigma_{Y,\mathrm{tr}}}.
\]
Per-seed values, means, standard deviations, and paired differences to each
principal control are retained.  Aggregate ranks or win--tie--loss summaries
are never shown without the underlying per-task uncertainty.

\paragraph{Result provenance.}
Every table and figure is generated from stored artifacts containing the
resolved generator, functional, split, model, candidate grid, predictions, and
metrics.  Failed or incomplete runs remain visible in the audit ledger, and
the supplementary material maps each reported number to its script and
artifact path.

\begin{table}[t]
\centering
\small
\caption{Staged empirical design.  The full Cartesian product of generators,
tasks, resources, and shot counts is intentionally avoided; each stage isolates
a distinct scientific question.}
\label{tab:empirical-stages}
\begin{tabularx}{\linewidth}{@{}lXX@{}}
\toprule
Stage & Main setting & Primary purpose \\
\midrule
E1 & Reference VARMA and functional hierarchy & Capacity, learning curves, and diagnostic task difficulty \\
E2 & VARMA, HMM, and GARCH across persistence levels & Memory--dependence interaction and cross-family robustness \\
E3 & High-dimensional sparse and distributed signals & JL distortion, fixed-width behavior, and dimensional failure modes \\
E4 & Principal classical controls & Representation-level comparative calibration \\
E5 & One-factor architecture sweeps & Accuracy--resource Pareto structure \\
E6 & Exact and finite-shot features & Measurement sensitivity and noisy-feature training \\
E7 & Lorenz--63 and Mackey--Glass & Out-of-scope chaotic stress testing \\
E8 & Three TSER datasets & Limited real-world transfer sanity check \\
\bottomrule
\end{tabularx}
\end{table}

\subsection{Experiment E1: Functional hierarchy and learning behavior}
\label{sec:exp-functional-hierarchy}

E1 uses the medium-persistence \(d=3\) VARMA realization.  With one fixed exact
featurizer, a regularization path records training MSE, validation MSE, test
MSE, RKHS norm, and Gram-matrix condition number for the future-projection,
exponential-memory, and Volterra tasks.  A nested-sample study uses
\[
  N_{\mathrm{tr}}\in\{100,500,1000,2000,3000,5000\}
\]
with the candidate grid fixed in advance.  Its purpose is to characterize
learning behavior, not to fit a prescribed \(N^{-1/2}\) law.  Delayed recall at
\(\tau\in\{1,5,10,20\}\) and the sparse cross-lag target
\((\tau_1,\tau_2)=(0,15)\) then isolate long-memory and nonlinear interaction
requirements.  The full target definitions are in
\cref{app:shared-functionals}.

\subsection{Experiment E2: Memory, persistence, and dependent-family
robustness}
\label{sec:exp-memory-dependence}

E2 varies the memory requirement of both the data and the target.  For VARMA,
HMM, and GARCH, the low-, medium-, and high-persistence realizations in
\cref{tab:dependent-family-suite} are paired with delayed recall and with
exponential-memory targets using a predeclared decay grid.  Around the selected
architecture, the contraction parameter is varied over
\[
  \lambda\in\{0.05,0.1,0.3,0.5,0.8,0.95\}.
\]
The main diagnostic is whether the contraction value minimizing validation
error changes systematically with process persistence or target delay, and
whether the relationship is non-monotonic.

A complementary cross-family comparison fixes the selected exact-feature
configuration and evaluates the exponential-memory and Volterra targets on all
three generators.  It additionally uses one-step future projection for VARMA,
the filtered-regime score for HMM, and the next conditional-volatility state
for GARCH.  This stage reports both absolute error and paired differences to the
principal classical controls; it is not a repeated architecture search for
each family.

\subsection{Experiment E3: Ambient dimension and JL width control}
\label{sec:exp-high-dimensional}

E3 uses the stable VARMA construction with
\(d\in\{3,10,30,100,300,500\}\) under the sparse and distributed target
directions in \cref{app:high-dimensional-functionals}.  The main curve fixes
\(n=5\), \(R=3\), \(k=2\), and the selected contraction while \(d\) grows.  A
width-sensitivity study uses \(n\in\{3,5,8\}\) at \(d\in\{100,500\}\).

For every run, we report prediction error, encoded width, feature dimension,
projection time, and the empirical distortion statistic
\begin{equation}
    \widehat\Delta_{0.95}
    :=
    \operatorname{quantile}_{0.95}
    \left(
      \left\{
      \left|
      \frac{\norm{\Pi x-\Pi x'}_2^2}{\norm{x-x'}_2^2}-1
      \right|:
      (x,x')\in\mathcal Q
      \right\}
    \right),
    \label{eq:empirical-jl-distortion}
\end{equation}
where \(\mathcal Q\) is a predeclared sample of training-point pairs.  The
comparison includes fixed-width \textsc{QuaRK}, classical JL plus Mat\'ern KRR,
matched random features, and raw-window readouts when tractable.  Sparse targets
test resistance to nuisance coordinates; distributed targets test whether
useful geometry survives compression when signal spans all coordinates.

\subsection{Experiment E4: Matched classical calibration}
\label{sec:exp-matched-controls}

E4 consolidates the principal comparisons from E1--E3 under common splits,
feature-dimension reporting, readout grids, and stochastic seeds.  For each
scenario it reports \textsc{QuaRK}, ESN plus Mat\'ern KRR, matched random
features plus Mat\'ern KRR, and classical JL plus Mat\'ern KRR.  Raw-window
models are included as uncompressed references.  Feature-dimension matching,
identical-readout matching, candidate-count matching, runtime, and memory are
reported as separate quantities rather than compressed into a single
``matched budget.''  The goal is to identify scenario-dependent strengths and
failures, not to establish universal dominance.

\subsection{Experiment E5: Resource and memory trade-offs}
\label{sec:exp-resource-tradeoffs}

E5 returns to medium-persistence VARMA and performs one-factor sweeps around the
selected configuration:
\[
  n\in\{3,5,6,8,10\},\qquad
  R\in\{1,3,5,8\},\qquad
  k\in\{1,2,3\},
\]
together with the contraction grid from E2.  The Volterra task and delayed
recall at \(\tau\in\{10,20\}\) are used because they require nonlinear
processing and retained history.  Each panel reports predictive error with
peak logical width, feature dimension, reservoir executions, exact-feature
runtime, and the implied state-preparation and channel-application counts.  A
Pareto summary identifies non-dominated accuracy--resource operating points.

\subsection{Experiment E6: Finite-shot degradation}
\label{sec:exp-finite-shot}

E6 fixes the selected reservoir and evaluates snapshot counts
\[
    S\in\{100,500,1000,5000,10000\}
\]
on the exponential-memory, Volterra, and \(\tau=10\) delayed-recall tasks.
Each setting is repeated with independent shadow randomness.  Two protocols are
reported: a readout trained on exact features and evaluated on estimated
features, and a readout trained and evaluated on independently estimated
features.  The exact feature vector is the reference.  In addition to task
error, we report coordinate-wise feature RMSE, Euclidean feature error,
prediction error, \(NRS\) state preparations, and \(NRSw\) temporal channel
applications.  Only the frozen-readout protocol is interpreted through
\cref{cor:measurement-stability}.

\subsection{Experiment E7: Out-of-scope chaotic stress tests}
\label{sec:exp-chaotic}

E7 evaluates Lorenz--63 and Mackey--Glass trajectories under the generators and
sampling rules in \cref{app:chaotic-stress}.  Forecast horizons, sampling
intervals, and observation-noise levels are varied while keeping chronological
splits and predeclared candidate grids.  \textsc{QuaRK} is compared with the
ESN and matched-random-feature controls using the same readout.  Forecast error
is reported as a function of horizon, together with the horizon at which each
method loses a predeclared fraction of the persistence-baseline skill.  These
experiments are empirical stress tests only: deterministic chaos is not
presented as satisfying the stationary beta-mixing assumptions of
\cref{thm:generalization}.

\subsection{Experiment E8: Limited real-world sanity checks}
\label{sec:exp-real-world-sanity}

E8 uses three heterogeneous datasets from the Monash--UEA--UCR time-series
extrinsic-regression archive \citep{TanEtAl2021TSER}: Gas Sensor Array Ethanol,
Live Fuel Moisture Content, and Hydraulic Systems.  They respectively provide
a very long univariate sensor trace, a year-long seven-band remote-sensing
series, and a short but very high-dimensional multivariate hydraulic record.
The audited archive splits, dimensions, preprocessing, and dataset roles are
specified in \cref{app:real-world-sanity}.

The comparison is deliberately small: \textsc{QuaRK}, ESN plus Mat\'ern KRR,
matched random nonlinear features plus Mat\'ern KRR, and raw-window Mat\'ern KRR
when feasible.  No stationarity or mixing claim is made for these datasets, and
three datasets are not treated as evidence of broad benchmark superiority.
Their purpose is to check whether failure modes and relative trends observed in
the controlled study persist in naturally collected data.  Full dataset-level
results belong in the appendix; the main text contains only a compact summary
with uncertainty and no aggregate claim beyond this limited transfer check.

\begin{center}
\fcolorbox{black}{gray!10}{
\begin{minipage}{0.90\linewidth}
\textbf{Draft status.}
The empirical objects and protocols above are fixed before the numerical tables
are inserted.  The final revision will populate E1--E8 from audited artifacts
and will update the abstract, discussion, and conclusion only with conclusions
supported by those artifacts.
\end{minipage}
}
\end{center}
